\documentclass[a4paper, oneside]{article}
\usepackage{graphicx}
\usepackage{amsthm}
\usepackage{amsmath}
\usepackage{amssymb}
\usepackage[a4paper,
            bindingoffset=0.2in,
            left=2cm,
            right=2cm,
            top=2cm,
            bottom=2cm,
            footskip=.25in]{geometry}
\usepackage[italian]{babel}
\usepackage{pgfplots}
\usepackage{tabularx}
\usepackage{tikz}
\usepackage{wrapfig}
\usepackage{color}
\usepackage[d]{esvect}
\definecolor{page}{rgb}{0.129,0.157,0.212}
\pagecolor{page}
\color{white}
\graphicspath{ {./images/} }
\usetikzlibrary{shapes.geometric}
\usetikzlibrary{datavisualization}
\usetikzlibrary{datavisualization.formats.functions}
\usetikzlibrary{patterns}
\pgfplotsset{width=10cm,compat=1.18}

\title{Appunti di Metodi}
\author{Tommaso Miliani}
\date{20-03-26}

\begin{document}
\newtheoremstyle{theoremEnv}
                {}          % Space above
                {}          % Space below
                {\slshape}  % Body font
                {}          % Indent amount
                {\bfseries} % Head font
                {.}         % Punctuation after head
                {\newline}  % Space after theorem head
                {}          % Theorem head spec
\theoremstyle{theoremEnv}

\newtheorem{definition}{Definizione}[section]
\newtheorem{theorem}{Teorema}[section]
\newtheorem{lemma}{Proposizione}[section]
\newtheorem{observation}{Osservazione}[section]
\newtheorem{corollary}{Corollario}[theorem]
\newtheorem{example}{Esempio}[section]
\newtheorem{remark}{Enunciato}[section]

\maketitle

\section{Dimostrazione calcolo di residui}
\begin{enumerate}
    \item Si assume che $f$ abbia un polo semplice in $a$, dunque il suo sviluppo di Laurent 
    vicino ad $a$ sarà della forma 
    \begin{gather*}
        f(z) = \frac{c_{-1}}{z - a} + c_0 + c_1(z - a) + \dots
    \end{gather*}
    Se si moltiplica per $z - a$:
    \begin{gather*}
        (z - a)f(z) = c_{-1} + c_0(z - a)+ c_{1}(z - a )^{2} + \dots
    \end{gather*}
    \item Il residuo è uguale al 
    \begin{gather*}
        \text{Res} = \lim_{z \to a} (z - a)\frac{h(z)}{g(z)} 
    \end{gather*}
    Poiché $g(a) = 0$, allora 
    \begin{gather*}
        \lim_{z \to a} (z - a)\frac{h(z)}{g(z) - g(a)} 
    \end{gather*}
    che è uguale a 
    \begin{gather*}
        \lim_{z \to a} \frac{h(z)}{\frac{g(z) - g(a)}{z - a}} = \frac{h(a)}{g'(a)} 
    \end{gather*}
    \item In $a$ la funzione $f$ ha un polo di ordine $n$. Dunque il suo sviluppo di Taylor 
    vicino ad $a$ ha la seguente forma (è regolare)
    \begin{gather*}
        f = \frac{1}{(z - a)^{m}} \sum_{n = 0}^{\infty } \frac{h^(n)(a)}{n!}(z - a )^{n} 
    \end{gather*}
    Dunque si ottiene lo sviluppo di Laurent: 
    \begin{gather*}
        = \sum_{n = 0}^{\infty } \frac{h^{(n)}(a)}{n!}(z - a)^{n - m}
    \end{gather*}
    si prende il residuo ed il coefficiente $n - m= - 1$ e $ n = m -1$. dunque è dimostrata.
    \item La dimostrazione parte dal caso 
    \begin{gather*}
        h = (z - a)^{m}f(z) \qquad f = \frac{(z - a)^{n}f(z)}{(z - a)^{m}} 
    \end{gather*}
    Applicando la terza, si trova la tesi.
\end{enumerate}

\section{Utilizzo dei residui per il calcolo degli integrali}
Si vuole ora calcolare alcuni integrali deifiniti. La tipologia dell'integrale non 
è assulutamente uno standard. 

\subsection{Primo tipo}
Sono degli integrali del tipo
\begin{align}
    I = \int_{0}^{2\pi} Q(\cos\theta, \sin\theta) \ d\theta 
\end{align}
questa tipologia di integrale richiede che la funzione $Q$ 
sia una funzione che coinvolga seno e coseno e che sia integrabile 
per un multiplo del periodo. È anche importante che $Q$ sia una funzione 
razionale (rapporto di due polinomi), definita sul cerchio
$\mathbb{S} = \{(x, y) : x^{2} + y^{2} = 1\}$. Seno e coseno parametrizzano i punti sul cerchio 
unitario, analogamente parametrizzati da $e^{i\theta}$; si può dunque 
eseguire un cambio di variabile complesso : 
\begin{gather*}
    z = e^{i\theta} = \gamma(\theta) \in \mathbb{C}
\end{gather*}
utilizzando $\gamma$ come contorno di integrazione al variare 
di $\theta \in [0, 2\pi]$, dunque $\gamma : [0, 2\pi] \to \mathbb{C}$ è un circuito. 
Si può ora vedere come sostituire $\cos\theta$ e $\sin\theta$:
\begin{align*}
    \cos\theta &= \frac{e^{i\theta} + e^{i\theta}}{2} = \frac{z + \frac{1}{z}}{2} \\
    \sin\theta &= \frac{e^{i\theta} - e^{-i\theta}}{2i} = \frac{z - \frac{1}{z}}{2i} \\
    dz &= ie^{i\theta} \ d\theta  = iz \ d\theta \ \Longrightarrow \ d\theta = \frac{dz}{iz}
\end{align*}
Dunque si può riscrivere 
\begin{align}
    I = \oint _\gamma Q\left(\frac{z + \frac{1}{z}}{2}, \frac{z - \frac{1}{z}}{2i}\right) \frac{dz}{iz} = \oint_{\left| z \right| = 1 } \tilde{Q}(z) \ dz
\end{align}
Applicando il teorema dei residui, l'integrale è determinato dalla somma dei residui dei poli di 
$\tilde{Q}$ all'interno del cerchio unitario, dunque 
\begin{gather*}
    I = 2\pi i \sum_{a \in Z(\tilde{Q}), \left| a \right| < 1 } \text{Res}(\tilde{Q}, a)
\end{gather*}
Si nota che $\left| a \right| <1$ in quanto non ci possono essere singolarità sulla circonferenza, 
in quanto corrisponderebbe alla singolarità della $Q$ normale 
sul circuito di integrazione ma, dato che essa è regolare, allora il circuito di 
integrazione finale non deve avere singolarità altrimenti non sarebbe ben definito. 

\begin{example}
    \begin{gather*}
        I = \int_{0}^{2\pi} \frac{1}{4\cos\theta - 5} \ d\theta  
    \end{gather*}
    Utilizzando le trasformazioni riportate sopra possiamo riscrivere 
    \begin{gather*}
        \oint _{\gamma} \frac{1}{4\frac{z + \frac{1}{z}}{2} - 5} \frac{dz}{iz} = \oint_{\gamma} \frac{-i}{z(2(z + \frac{1}{z}) - 5)} \ dz = \oint_{\gamma} \frac{-i}{2z^{2} - 5z + 2} \ dz = \oint_{\gamma} \frac{-i}{2(z - 2)(z - \frac{1}{z})}dz
    \end{gather*}
    Dunque si cercano i poli di questa funzione: ossia $z = 2$ e $z = \frac{1}{2}$. Solo il secondo è 
    all'interno del cerchio unitario, dunque considero solo quel polo. Allora 
    il residuo, calcolato in $\frac{1}{2}$ è 
    \begin{gather*}
        \text{Res} = - \frac{i}{2(\frac{1}{2} - 2)} = \frac{i}{3}
    \end{gather*}
    Dunque
    \begin{gather*}
        I = 2\pi i \frac{i}{3} = - \frac{2\pi}{3}
    \end{gather*}
\end{example}

\subsection{Secondo tipo}
Sono integrali del tipo 
\begin{align}
    I = \int_{\mathbb{R}} f(x) \ dx 
\end{align}
dove $f : \mathbb{R} \to \mathbb{C}$ tale che può essere estesa 
al piano complesso ad una funzione olomorfa (ed integrabile). Basta che 
sia olomorfa su almeno uno dei due semipiani $Im(z) \geq 0$ o $Im(z) \leq 0$. Per 
olomorfa si intende di \textbf{almeno di singolarità isolate}, chiaramente queste singolarità 
isolate devono essere al di fuori dell'asse reale, altrimenti non sarebbe integrabile 
sui reali ($Im(a) \neq 0$). Se esiste 
\begin{gather*}
    \lim_{r \to \infty } \int_{-r}^{r} f(x)  \ dx  
\end{gather*}
in quanto la funzione, anche se non è integrabile, è integrabile con questo 
limite (un esempio è il seno di x, anche se non è del secondo tipo). In questo 
caso si può dunque integrare sull'asse reale. Per ottenere un percorso 
chiuso si può aggiungere una superficie di raggio $r$ con $r \to \infty $. Questo perché 
sull'asse reale si ottiene l'integrale appena scritto, mentre sulla semicirconferenza 
devo calcolare un integrale differente che, se tende a zero, allora l'integrale sul circuito 
corrisponde solo all'integrale sull'asse reale $I$. Supponendo che
$f$ sia un rapporto tra polinomi $\frac{p(x)}{q(x)}$, in cui $q(x) \neq 0, \forall x \in \mathbb{R}$, 
l'integrale esiste se $\deg(q) \leq \deg(p) + 2$ poiché $\lim_{x \to \infty }xf(x)  = 0$, dunque devo avere 
un andamento simile a $\frac{1}{x^{\geq 2}}$ per poter convergere. Dunque 
\begin{gather*}
    I = \int_{-r}^{r} f(x)  \ dx = \int_{\beta_r} f(z) \ dz
\end{gather*}
Mentre l'integrale sul percorso completo 
\begin{gather*}
    \oint_{\beta_r + \gamma_r} f(z) \ dz = \int_{\beta_r} f(z) \ dz + \int_{\gamma_r} f(z) \ dz 
\end{gather*}
Il primo è esattamente $I$, mentre quello su $\gamma_r$ è l'integrale 
\begin{gather*}
    \int_{\gamma_r}^{}  \frac{p(z)}{q(z)} \ dz 
\end{gather*}
in cui $\gamma_r$ + parametrizzabile come $\gamma_r = re^{i\theta}, \theta \in [0, 2\pi]$ con $r \to \infty $:
\begin{gather*}
    \int \frac{p(re^{i\theta})}{q(re^{i\theta})} rie^{i\theta} \ d\theta 
\end{gather*}
Allora posti gli andamenti $p \sim x^{K}$ e $q \sim x^{L}$ e dato che $L \geq K + 2$, si ottiene che 
l'andamento dell'integrale è $\sim r^{K - L + 1}$ ossia $\sim r^{\leq 1}$, dunque convergerà. Allora 
\begin{gather*}
    I = \lim_{r \to \infty } \oint_{\beta + \gamma_r} f(z) \ dz = 2\pi i \sum_{Im(a) > 0 } \text{Res}(f, a)
\end{gather*}
Questo si può anche estendere alla semicirconferenza sotto l'asse reale che, in questo caso, 
l'integrale dovrà tenere un segno meno per ricordare che si percorre la curva al contrario, ottenendo
\begin{gather*}
    -2\pi i \sum_{Im a < 0 } \text{Res}(f, a) 
\end{gather*}

\begin{lemma}[Lemma del cerchio grande]
    Presi due angoli $\theta_1 < \theta_2$ con $\theta_2 - \theta_1 \leq 2\pi$, che individuano 
    il settore circolare 
    \begin{gather*}
        A(\theta_1, \theta_2) \equiv \{z \in \mathbb{C} : \ \theta_1 \leq \arg z \leq \theta_2\}
    \end{gather*}
    In questo settore circolare si definisce la funzione $f : A(\theta_1, \theta_2) \to \mathbb{C}$, con singolarità isolate 
    (al più e numero finito). Preso $C_r$ l'arco di circonferenza di raggio $r$ in quel settore
    \begin{gather*}
        C_r = \{z \in \mathbb{C} : \ \left| z \right| = r, \theta_1 \leq \arg(z) \leq \theta_2 \}
    \end{gather*}
    Se $zf(z)$ tende a zero \textbf{uniformemente} in $A(\theta_1, \theta_2)$ per $\left| z \right| \to \infty$, ovvero se
    \begin{gather*}
        \lim_{r \to \infty } \sup_{z \in C_r} \left| z f(z)\right|  = 0
    \end{gather*}  
    allora il limite 
    \begin{gather*}
        \lim_{r \to \infty } \int_{C_r}^{} f = 0  
    \end{gather*}
\end{lemma}
\begin{proof}
    Posto
    \begin{gather*}
        M(r) = \sup_{z \in C_r} \left| f(z) \right|  
    \end{gather*}
    l'estreo superiore di $|f|$ su $C_r$, dunque 
    \begin{gather*}
        \left| \int_{C_r} f\right|   \leq L(C_r)M(r) = \left| \theta_2 - \theta_1 \right|rM(r ) 
    \end{gather*}
    Per $r \to \infty$, $rM(r) \to 0$ per la condizione che $\left| zf(z) \right| \to 0$.  
\end{proof}

\begin{example}
    Si vuole integrare
    \begin{gather*}
        I = \int_{\mathbb{R}} \frac{1}{x^{2} + 1} \ dx 
    \end{gather*}
    dove i poli sono a $\pm i$. 
\end{example}

\section{Terzo tipo}
Sono anch'essi integrali su tutto $\mathbb{R}$ della forma
\begin{align}
    \int_{\mathbb{R}} f(x) e^{i\alpha x} \ dx 
\end{align}
ma con la differenza che c'è un esponenziale.
\begin{itemize}
    \item Nel caso in cui $\alpha = 0$, 
ci si riconduce agli integrali di secondo tipo. 
    \item $\alpha > 0$ si ha l'integrale sul circuito 
    nel semipiano immaginario positivo e si fa utilizzo del \textbf{Lemma di Jordan}.
    \item $\alpha < 0$ si ha l'integrale sul circuito nel 
    semipiano immaginario negativo, riconducendosi al caso $\alpha > 0$,
\end{itemize}

\begin{lemma}[Lemma di Jordan]
    Sia $f: A(\theta_1, \theta_2) \to \mathbb{C}$ una funzione definita in un settore circolare
    tale che $0 \leq \theta_1 < \theta_2 \leq \pi$. Sia $\gamma_r :[\theta_1, \theta_2] \to \mathbb{C}$, 
    $\gamma_r(\theta) = re^{i\theta}$ l'arco di circonferenza di 
    raggio $r > 0 $ in $A(\theta_1, \theta_2)$. Se 
    \begin{gather*}
        \lim_{z \to \infty } f(z) = 0 
    \end{gather*}
    uniformemente rispetto all'angolo e $\alpha > 0$ allora 
    \begin{align}
        \lim_{r \to \infty } \int_{\gamma_r} f(z) e^{i\alpha z} \ dz = 0  
    \end{align}
\end{lemma}
\begin{proof}
    Non richiesta.
\end{proof}

Se le ipotesi del teorema sono soddisfatte, allora si possono calcolare 
gli integrali come i seguenti
\begin{itemize}
    \item $\alpha > 0$:
    \begin{gather*}
        \int_{\mathbb{R}} e^{i\alpha x} f(x) \ d x= 2\pi i \sum_{a : Im(a) > 0 } \text{Res}(fe^{i\alpha z}, a)  
    \end{gather*}
    \item $\alpha < 0$:
    \begin{gather*}
        \int_{\mathbb{R}} e^{i\alpha x} f(x) \ d x = -2\pi i \sum_{a : Im(a) < 0 } \text{Res}(fe^{i\alpha z}, a)
    \end{gather*}
\end{itemize}
Si osserva come il risultato sia lo stesso, ma ricordandosi di cambiare segno. Il Lemma di Jordan 
è meno restrittivo del teorema dei cerchi grandi.

\begin{example}[Dispense]
    
\end{example}

\end{document}